The AIND and the SWC/GCNU share a complementary vision for advancing systems neuroscience. Both institutions are committed to cutting-edge experimental approaches, computational modeling, open science, and technology development. A strong spirit of collaboration and mutual respect exists between our organizations, built on shared scientific priorities and ongoing interactions.

\paragraph{Key points of connection}

\begin{description}

\item[Shared goals and open science ethos:]  
Both the Allen Institute and SWC/GCNU emphasize the importance of generating high-quality, large-scale neuroscience datasets and making them openly available to the scientific community. This shared commitment to openness and reproducibility fosters natural synergies.

\item[Complementary expertise:]  
The Allen Institute brings deep expertise in large-scale neural recordings, cloud-based infrastructure, and advanced visualization platforms. The SWC/GCNU contributes the AEON platform, several years of experience with NaLoDuCo experimentation, and strong capabilities in computational neuroscience and statistical modeling.

\item[Scientific advising and leadership:]  
Prof.~Peter Dayan (former director of the GCNU) and Prof.~Thomas Mrsic-Flogel (current director of the SWC) both serve as scientific advisors to the AIND. In addition, Dr.~Gary Wilson, Head of Science Portfolio at the Gatsby Charitable Foundation, advises the Allen Institute's Frontiers Group.

\item[Development of Neuropixels:]  
The Neuropixels probe was developed through a consortium of funders including the Howard Hughes Medical Institute (HHMI), the Wellcome Trust, the Gatsby Charitable Foundation, and the Allen Institute for Brain Science. The Wellcome Trust and the Gatsby Foundation are the primary funders of the SWC, and SWC scientists played key roles in the development of this technology.

\item[Ongoing collaborations:]  

Scientists from SWC/GCNU and the Allen Institute interact regularly at
        scientific meetings and symposia and through infrastructure discussions
        (e.g., the AEON platform). In 2021, Dr.~Christof Koch, Meritorious
        Investigator at the Allen Institute, delivered a talk at the SWC about
        the Allen Institute’s research culture. That same year, Dr.~Carl
        Schoonover (Senior Scientist at AIND and co-lead on this proposal)
        presented a talk on olfactory learning and forgetting. In 2023,
        Dr.~Adam Tyson (Head of Neuroinformatics at SWC/GCNU) visited AIND to
        exchange best software practices with the Scientific Computing team.

\end{description}

\subsubsection*{Collaboration on NaLoDuCo Experimentation}

Both the SWC/GCNU and AIND have a longstanding interest in the neural basis of food foraging behavior.

In 2020, the SWC and GCNU launched the Foraging Working Group—uniting all experimental and computational research groups across the centers—in an unprecedented team-science effort. This led to the development of the AEON platform and the initiation of studies into the neural computations underlying NaLoDuCo foraging behavior in mice.

In 2021, the newly founded AIND made foraging a central pillar of its research
program. While initial efforts focused on head-fixed virtual reality setups, in
2023, Dr.~Carl Schoonover joined AIND to launch a new NaLoDuCo research
initiative focused on odour learning in freely moving animals using the AEON platform.

This joint grant proposal marks the beginning of our formal collaboration on
foraging research. We are strongly aligned in both vision and scientific
objectives, and our respective programs bring highly complementary expertise.
By co-developing shared infrastructure for NaLoDuCo experimentation, we aim to
accelerate progress in this field and provide broadly accessible tools for the
global neuroscience community.

\subsubsection*{Planned Collaborative Activities}

This collaboration will substantially strengthen the partnership between two world-leading neuroscience institutions. Planned collaborative activities include:

\begin{itemize}
  \item \textbf{Monthly virtual coordination meetings} between SWC/GCNU and AIND investigators to align development goals, share progress, and troubleshoot challenges.
  
  \item \textbf{Annual in-person workshops}, alternating between the US and the UK, to assess progress against project milestones and refine the collaborative roadmap.
  
  \item \textbf{Joint journal clubs}, held monthly, to discuss recent research relevant to NaLoDuCo experimentation and analysis.
  
  \item \textbf{Cross-institutional benchmarking} of software tools using representative datasets from both sites, ensuring broad applicability across experimental platforms and analysis pipelines.
  
  \item \textbf{Shared software dissemination} through collaborative repositories that reflect the joint development and commitment of both institutions.
  
  \item \textbf{Co-authored publications} highlighting methodological innovations and scientific insights gained from long-duration, continual data.
  
  \item \textbf{Community outreach}, including hands-on tutorials, conference workshops, and training sessions to promote NaLoDuCo methods and foster adoption in the broader neuroscience community.
\end{itemize}

